The definitions given in this section are a direct adaptation from \cite{JOHANSSON20061}, \cite{SERFOZO19901}, \cite{ManuelaGirottiPhD}. 

\section{Point Processes}

 In general, point processes are models to describe random numbers of occurrences of events in a given time-frame or, alternatively, the number of points in regions of some space. For this work, we mainly focus on the latter application: density of points in regions of space. To illustrate such processes, suppose there are $m$ points at the random locations $\chi_m = \{\mmany{x}{m}\}$ in some region $\Lambda$ of the Euclidean space $\R^n$. We can define a quantity $\xi(A)$ for any sub-region of $A \subset \Lambda$ that measures the amount of points ($\#\{\chi_m \cap A\}$) in this region in the following way:
\begin{equation}
	\xi(A) = \sum_{k=1}^{m} \mathbb{1}(x_k \in A).
\end{equation}
Let $\mathcal{I}$ denote all intervals or rectangles in $\Lambda$, $\mathcal{E}$ denote the class of Borel sets and $\mathcal{B}$ the class of bounded Borel sets. One can notice that this $\xi$ defines a \textit{random counting measure}, i.e., a mapping $\mu: \mathcal{E} \rightarrow \{0,1,\cdots, \infty\}$ such that 
\begin{equation}
	\mu(\emptyset) = 0, \quad \mu(B) < \infty, \ \text{for} \ B \in \mathcal{B},
\end{equation} 
and, for disjoint $B_1, B_2, \cdots$ in $\mathcal{B},$
\begin{equation}
	\mu\left( \bigcup_k B_k \right) = \sum_k \mu(B_k).
\end{equation} 
We say that $\xi$ is boundedly finite. Let $\mathcal{N}(\Lambda)$ denote the space of all counting measures $\xi$ on $\Lambda$. If $B \in \mathcal{B}$, then $\xi(B)$ is finite and we can write
\begin{equation}
	\xi|_B = \sum_{i=1}^{\xi(B)} \delta_{x_i},
\end{equation}
for some $\chi = \{\mmany{x}{{\xi(B)}}\} \subset \Lambda$. Finally, we can explicitly define point processes.

\begin{definition} (\textit{Point Processes})
	Let $\Lambda$ be a complete separable metric space. A \textbf{point process} $\mathbb{P}$ on $\Lambda$ is a probability measure on the space of all configuration of points $\chi$ or counting measures $\mathcal{N}(\Lambda)$. We say $\mathbb{P}$ to be a $n$-point process if $\mathbb{P}(\# \chi = n) = 1$ or, equivalently, $\mathbb{P}(\xi(\Lambda) = n)  = 1$. We say that $\xi$ is simple if $\xi(\{x\}) \leq 1$ for all $x \in \Lambda$. The point process is simple if $\mathbb{P}( \xi \ \text{is simple}) = 1$.
\end{definition}

An important note is that $\xi$ is the same for any ordering of the subscripts on \many{x}{n}. This lead us the the conclusion that if $u_n(\mmany{x}{n})$ is a probability function on $\R^n$ such that it is invariant under permutations, i.e.
\begin{equation}
	u_n(x_{\sigma(1)}, \cdots, x_{\sigma(n)}) = u_n(\mmany{x}{n}), \quad \forall \sigma \in S_n,
\end{equation}
then $u_n$ defines naturally a point process.

The mapping $\mathcal{E} \ni A \mapsto \E[\#(\chi \cap A)]$, that assigns to the the Borel set $A$ the expected value of the number of points in itself under the configuration $\chi$, is a measure on $\R$. The mean measure is defined by 
\begin{equation}
	\mu(A) = \E[\xi(A)] = \E[\#(\chi \cap A)], \quad A \in \mathcal{E}.
\end{equation}

This $\mu$ is a measure on $\Lambda$ in the standard sense that $\mu : \mathcal \rightarrow [0, \infty]$ satisfies $\mu(\emptyset)) = 0$ and, for disjoint $\Lambda_1, \Lambda_2, \cdots$ in $\mathcal{E}$, 
\begin{equation}
	\mu\left( \bigcup_k \Lambda_k \right) = \sum_k \mu(\Lambda_k).
\end{equation} 
Assuming there exists a density $\rho_1: \mathcal{E} \ni A \rightarrow \R_+$ with respect to the Lebesgue measure (\textbf{$1$-point correlation function}), there is also a way to define the mean measure as 
\begin{equation}
	\mu(A) = \int_A \rho_1(x) \dd x
\end{equation}
This can be called the rate or intensity of $\xi$ and is the infinitesimal probability of a point at $x$, i.e. $\mathbb{P}(\xi(\dd x) = 1) = \rho_1(x) \dd x$. In this case $\mu$ is written as $\mu(\dd x) = \rho_1(x) \dd x$. This mean measure shows up in expressions of expectations for functionals of $\xi$. More generally, the \textbf{k-th moment} $\mu_k$ of $\xi$ , written
\begin{equation}
	\mu_k(\mcmany{A}{k}{\times}) = \E \left[ \xi_1(A) \xi_2(A)\cdots \xi_k(A) \right], \quad \mmany{A}{k} \cdots \in \mathcal{E},
\end{equation}
can be expressed in respect to the \textbf{k-th correlation function} $\rho_k: A^k \rightarrow \R_+$ in the following way
\begin{equation}
	\mu_k(\mcmany{A}{k}{\times}) = \E \left[ \prod_{j=1}^{k} (\# \chi \cap A_j) \right]= \int_{A_1} \cdots \int_{A_k} \rho_k(x_1, \cdots, x_k) \dd x_1 \cdots \dd x_k.
	\label{Equation: kth moment measure}
\end{equation}
That is, the expected number of $k$-tuples $\chi_k = \{\mmany{x}{k}\} \in \chi^k$ such that $x_j \in A_j\ \forall j.$ For example, if $u_n(\mmany{x}{n})$ is a probability density function on $\R^n$, invariant under permutation of coordinates, then we can build an $n$-point process on $\R$ with correlation functions
\begin{equation}
	\rho_k(\mmany{x}{k}) \deff \frac{n!}{(n-k)!} \int_{\R^{n-k}} u_n(\mmany{x}{n}) \dd x_{k+1} \cdots \dd x_n.
	\label{Equation: Symmetric correlation}
\end{equation}

The mean measure also arises in expressions of expectations of functionals of $\xi$. For example, the Campbell formula that, for $E \in \mathcal{E}$, states that the integral
\begin{equation}
	\int_E f(x) \xi(\dd x) = \sum_{k=1}^{m} f(x_k),
\end{equation}
where $f: E \rightarrow \R_+$, has expectation
\begin{equation}
	\E\left[ \int f(x) \xi(\dd x) \right] = \int f \mu(\dd x)
\end{equation}
if the integral $\int_E |f(x)| \mu(\dd x)$ is finite.  

For simple points processes on $\R$ there is a direct way to understand correlation functions. Let $A_i = [x_i, x_i + \Delta x_i], 1 \leq i \leq k$, be disjoint intervals. If the $\Delta x_i$ are small enough, then we should expect either one or zero particles in each interval $A_i$. Hence, the product $\xi(A_1) \cdots \xi(A_k)$ is typically either $1$ or $0$ depending on whether we have one particle in each interval. From \eqref{Equation: kth moment measure}, we should expect
\begin{equation}
	\rho_k(\mmany{\lambda}{k}) = \lim_{\Delta x_i \rightarrow 0} \frac{\mathbb{P}(\text{one particle in each} \ A_i)}{\mcmany{\Delta x}{k}{}}.
\end{equation}
This means that $\rho_1(\lambda)$ is the density of particles at $\lambda$, but since we have many particles, the events of finding particles at their respective points are not disjoint events, even for distinct points. It follows from the definition that if we have a simple point process, then $\rho_k(\mmany{x}{k})$ is exactly the probability of finding particles at \many{x}{n}.

Lets explore an important example for our work. Let $\mathcal{H}_N$ be the space of all $N\times N$ Hermitian matrices. This space can be identified naturally with $\R^{N^2}$. If $\mu_N$ is a probability measure on $\mathcal{H}_N$ and $\{\lambda_1(M), \cdots, \lambda_N(M) \}$ denotes the set of eigenvalues of $M \in \mathcal{H}_N$, then
\begin{equation}
	M \rightarrow \sum_{j=1}^{N} \delta_{\lambda_j(M)}
\end{equation}
maps $\mu_N$ to a finite point process $P$ on $\R$. If $\dd M$ is the Lebesgue measure on $\mathcal{H}_N$, then
\begin{equation}
	\dd \mu_N(M) = \frac{1}{\pZ{N}{}} \ee^{-\tr{M^2}} \dd M
\end{equation}
is a Gaussian probability measure on $\mathcal{H}_N$, called the $GUE$. It can be shown that for any symmetric, continuous function $f$ on $\R^N$ with compact support
\begin{equation}
	\int_{\mathcal{H}_N} f(\lambda_1(M), \cdots, \lambda_N(M)) \dd \mu_N(M) = \int_{\R^N} f(\mmany{x}{N}) u_N(\mmany{x}{N}) \dd^N x 
\end{equation}
where
\begin{equation}
	u_N(\mmany{x}{N}) = \frac{1}{\pZ{N}{} N!} \prod_{1 \leq i < j \leq j} (x_i - x_j)^2 \prod_{j=1}^{N} \ee^{-x^2_j}
\end{equation}
is the induced eigenvalue measure on $\R^N$. Hence the point process $P$ on $\R$ defined by this induced measure has correlation functions given by \eqref{Equation: Symmetric correlation}.

 Most, if not all, of the processes discussed on this work will have an rightmost or last particle. We say that a point process $\xi$ has a last particle if there is a $t \in \R$ such that $\xi(t, \infty) < \infty$. If $\mcmany{x}{{{n(\epsilon)}}}{\leq}$ are finitely many particles in $(t, \infty)$, we define $x^* = x_{n(\xi)}$, the position of the last particle. The distribution $\mathbb{P}[x^* \leq t]$ is called the last particle distribution.

\begin{prop}
	Take $\xi$ a point process on $\R$, or a subset of, all whose correlation functions $\rho_k$ exist, and assume that
	\begin{equation}
		\sum_{n=0}^{\infty} \frac{1}{n!} \int_{(t, \infty)^n} \rho_n(\mmany{x}{n}) \dd \mu_n(x) < \infty
		\label{Equation: premisse last particle}
	\end{equation}
	for each $t \in \R$. then the process $\xi$ has a last particle and 
	\begin{equation}
		\mathbb{P}[x^* \leq t] = \sum_{n=0}^{\infty} \frac{(-1)^n}{n!} \int_{(t,\infty)^n} \rho_n(\mmany{x}{n}) \dd \mu_n(x) .
		\label{Equation: conclusion last particle}
	\end{equation}
\end{prop}
\begin{proof}
	We take $t < s$, then
	\begin{equation}
		\eP[\text{no particle in} \ (t, s)] = \sum_{n=0}^{\infty} \frac{(-1)^n}{n!} \int_{(t, s)^n} \rho_n(\mmany{x}{n}) \dd \mu_x(x).
	\end{equation}
	We see from \eqref{Equation: premisse last particle} and the dominated convergence theorem that we can let $s \rightarrow \infty$ and obtain \eqref{Equation: conclusion last particle}.
\end{proof}

%If $\rho_N$ is a symmetric probability density on $\R^N$, then %$(\mmany{x}{N}) \rightarrow \sum_{i=1}^N \delta_{x_i}$ maps the %probability measure with density $\rho_N$ to a finite point process on %$\R$. The correlation functions are given by
%\begin{equation}
	%\mu_n(\mmany{x}{n}) = \frac{N!}{(N - n)!} \int_{\R^{N-n}} %\rho_n(\mmany{x}{N}) \dd x_{n+1} \cdots \dd x_N,
%\label{Equation: correlation hermitian}	
%\end{equation}
%i.e., they are multiple marginal densities. 

\section{Determinantal Point Processes}

Repulsive random processes are, in general, intractable, that it, they are usually difficult of even impossible to, for example, compute marginals and conditional probabilities. In these cases, their theoretical analysis is a major challenge. However, there is a class of processes, repulsive ones, that are an rare exception to this rule: Determinantal Point Processes (DDP), or Fermionic point processes in the physics literature, are random processes that have the property of being both repulsive and tractable \cite{Tremblay}. This property comes from the fact that their correlation function can be expressed as a certain determinantal form. Let's make this precise.

\begin{definition}
	Consider a point process $P$ on a complete separable metric space $\Lambda$, with reference measure $\mu$, all of whose correlation functions $\rho_k$ exist. If there is a function $K: \Lambda \times \Lambda \rightarrow \C$ such that 
	\begin{equation}
		\rho_k(\mmany{x}{k}) = \det(K(x_i, x_j))_{i,j = 1}^{k}
	\end{equation}
	for all $\mmany{x}{k} \in \Lambda, \ k \geq 1$, then we say that $P$ is a \textbf{determinantal point process}. We call $K$ the \textbf{correlation kernel} of the process. 
\end{definition}

Important examples of DDP's are constructed by using the following result

\begin{thm} (Construction for DDP \cite{ManuelaGirottiPhD})
	Consider a kernel $K: \R \times \R \rightarrow \R_+$ with the following properties:
	\begin{itemize}
		\item Trace-class: $\tr(K) = \int_{\R} K(x,x) \dd x = n < + \infty$;
		\item Positivity: $\det[K(x_i, x_j)]_{i,j=1}^n$ is non-negative for every $\mmany{x}{n} \in \R$;
		\item Reproducing property: $\forall x, y \in \R$
		\begin{equation}
			K(x,y) = \int_{R} K(x, s) K(s, y) \dd s.
		\end{equation}
	\end{itemize}
	Then, 
	\begin{equation}
		u_n(\mmany{x}{n}) \deff \frac{1}{n!} \det[K(x_i, x_j)]_{i,j = 1}^n
	\end{equation}
	is a probability measure on $\R^n$, invariant under coordinated permutations. The associated point process is a determinantal point process with $K$ as a correlation kernel.
\end{thm}

All important information is contained in the correlation kernel for these processes and, therefore, all quantities of interest can be expresses i terms of the same $K$. The kernel $K$ can be seen as an integral kernel of an operator $K$ on $L^2(\Lambda, \lambda)$,
\begin{equation}
	Kf(x) = \int_\Lambda K(x, y) f(y) \dd \lambda(y)
	\label{Equation: integral operator K}
\end{equation}
as long as the right and side is well defined.

\subsection{Gap probabilities}

Consider a point process on $\R$ with correlation function $\rho_k$ and let $A \in \mathcal{E}$ be a Borel set such that, with probability $1$, there are only finitely may points in $A$ (for example, $A \in \mathcal{B}$). Denote by $\rho_A(n)$ the probability that there are exactly $n$ points in $A$. If there are $n$ points in $A$, then the number of ordered $k$-tuples in $A$ is $\binom{n}{k}$. Also, remember by \eqref{Equation: kth moment measure}, that 
\begin{equation}
\E \left[ \prod_{j=1}^{k} (\# \chi \cap A_j) \right]= \frac{1}{k!} \int_{A} \cdots \int_{A} \rho_k(x_1, \cdots, x_k) \dd x_1 \cdots \dd x_k
\end{equation}
should be the expected number of ordered $k$-tuples $(\mmany{x}{k})$ such that \cmany{x}{k}{<} and $x_j \in A$ for every $j = 1, \cdots, k.$ Therefore it must hold that
\begin{equation}
	\frac{1}{k!} \int_{A^k} \rho_k(\mmany{x}{k}) \dd x_1 \cdots \dd x_k = \sum_{n=k}^\infty \binom{n}{k} \rho_A(n).
\end{equation}
Assume the following alternating series is absolutely convergent, then
\begin{align}
	\sum_{k=0}^{\infty} \frac{(-1)^k}{k! } \int_{A^k} \rho_k(\mmany{x}{k}) \dd x_1 \cdots \dd x_k & = \sum_{k=0}^{\infty} \sum_{n=k}^\infty (-1)^k \binom{n}{k} \rho_A(n), \\
	& = \sum_{n=0}^\infty \left( \sum_{k=0}^{n} (-1)^k \binom{n}{k} \right) \rho_A(n).
\end{align}
On the other hand, $\sum_{k=0}^\infty (-1)^k \binom{n}{k}$ vanishes unless $n$ is zero. In conclusion,
\begin{equation}
	\rho_A(0) = \sum_{k=0}^{\infty} \frac{(-1)^k}{k! } \int_{A^k} \rho_k(\mmany{x}{k}) \dd x_1 \cdots \dd x_k,
\end{equation}
where we call $\rho_A(0)$ the \textbf{gap probability}, i.e. the probability of finding $0$ particles in $A$. In particular, if the point process is determinantal, we have
\begin{equation}
	\rho_A(0) = \sum_{k=0}^{\infty}  \frac{(-1)^k}{k! } \int_{A^k} \det[K(x_i, x_j)]_{i,j = 1}^k \dd x_1 \cdots \dd x_k.
\end{equation}
Let $d(z) = \det(I + zA)$ be the Fredholm determinant given by
\begin{equation}
	d(z) = \sum_{k=0}^{\infty}  \frac{z^k}{k! } \int_{[a,b]^k} \det[K(x_i, x_j)]_{i,j = 1}^k \dd x_1 \cdots \dd x_k.
\end{equation}
then it is clear that $\rho_0(A)$ is the Fredholm determinant $d(-1) = \det(I - \bold{K}|_A)$ of the integral operator
\begin{align}
	\bold{K}|_A &: L^2(A) \rightarrow L^2(A) \\
	& f(x) \mapsto \bold{K}[f](x) = \int_{\R} K(x,y) f(y) \dd y.
\end{align}
as defined in \eqref{Equation: integral operator K}. A more general result can be posed as follows:

\begin{thm} \cite{ManuelaGirottiPhD}
	Consider a determinantal point process with kernel $K$. For any finite Borel sets $A_j \subset \R, j = 1, \cdots, n$, the generating function of the probability distribution of the occupation number $\#_{A_j} \deff \# \{x_i \in A_j\}$ is given by
	\begin{equation}
		G(\mmany{z}{n}) = \E \left( \prod_{i=1}^{n} z_j^{\#_{A_j}} \right) = \det \left( I + \sum_{j=1}^n (z_j - 1) K|_{A_j} \right)
	\end{equation}
\end{thm}

From the properties of generating functions, i.e. $G(z) = \mathbb{E}(z^X) = \sum_{x=0}^\infty p(x) z^x$, it is true that
\begin{equation}
	p(k) = \mathbb{P}(X = k) = \frac{G^{(k)}(0)}{k!}.
\end{equation}
 Then, for example,  the probability of having $k$ points in $A$ is
 \begin{equation}
 	\mathbb{P}(\chi \cap A = k) = \frac{G^{(k)}(0)}{k!} = \frac{1}{k!} \frac{\dd^k}{\dd z^k} \det \left( I + (z-1) K|_A\right)_{z=0}.
 \end{equation}
 
\subsection{Outermost Particle}

\begin{prop}
	Consider a determinantal point process $\xi$ on a subset $\Lambda$  of $\R$ with a hermitian correlation kernel $K(x,y)$, i.e., $K(y,x)  = \bar{K(x, y)}$. Assume that $K(x,y)$ defines a trace class operator $K$ on $L^2(t, \infty)$ for each $t \in \R$, and that
	\begin{equation}
		\tr K = \int_t^\infty K(x,x) \dd \lambda (x) < \infty.
	\end{equation}
	Then $\xi$ has a last particle almost surely and 
	\begin{equation}
		\eP[\lambda^* \leq t] = \det{(I - K)}_{L^2(t, \infty)}
	\end{equation}
\end{prop}
\begin{proof}
	Since correlation functions are non negative and $K(y,x)  = \bar{K(x, y)}$, the matrix $K(x_i, x_j)$ is positive definite. In thta case Hadamard's inequality says that 
	\begin{equation}
		\det\left( K(x_i, x_j) \right)_{1 \leq i, j \leq n} \leq \prod_{j=1}^{n} K(x_j, x_j).
	\end{equation}
	Hence,
	\begin{align}
		\sum_{n=0}^\infty \frac{C^n}{n!} \int_{(t, \infty)^n} \det(K(x_i, x_j))_{1 \leq i, j \leq n} \dd^n \lambda(x)& \leq \sum_{n=0}^{\infty} \frac{C^n}{n!} \left( \int_t^\infty K(x,x) \dd \lambda(x) \right) \\
		& \exp{\left(C \int_t^\infty K(x,x) \dd \lambda(x) \right)}.
	\end{align}
\end{proof}
	
Consider for example the \textit{Airy Kernel},
\begin{equation}
	A(x, y) = \int_{0}^{\infty} \Ai(x+t) \Ai(y + t) \dd t
\end{equation}
This kernel is Hermitian and has 
\begin{equation}
	\int_{t}^{\infty} A(x,x) \dd x < \infty, \quad \forall t \in \R.
\end{equation}
It can be shown that $A(x,y) \_{(t,\infty)}$ is the kernel of a trace class operator and there is a point process $\xi$ on $\R$, the \textit{Airy kernel point process}, which is determinantal with kernel $A(x,y)$. We have
\begin{equation}
	F_{TW}(t) \deff \eP[\lambda^*(\xi) \leq t] = \det(I - A)_{L^2(t, \infty)}.
\end{equation} 
The distribution function $F_{TW}(t)$ for the last particle in the Airy kernel point process is a natural scaling limit of certain finite determinantal point processes. We call it the \textit{Tracy-Widow distribution}.

\subsection{Biorthogonal Ensemble}

Let $p_t(x; y)$ be the transition probability density from point $x$ to point $y$ at time $t$ of a one-dimensional strong Markov process with continuous sample paths. The following infamous theorem give us a determinantal formula for the probability that a number of paths with given starting and ending positions fall in certain sets at some time without intersecting in the mean time.

\begin{thm}
(Karlin, McGregor \cite{Karlin-McGregor}) Consider $n$ independent copies $X_1(t), \cdots , X_n(t)$ of a one-dimensional strong Markov process with continuous sample paths, conditioned so that
\begin{equation}
	X_j(0) = a_j, \quad \mcmany{a}{n}{<} \in \R.
\end{equation}
Let $p_t(x, y)$ be the transition probability function
of the Markov process and let $E_1,\cdots, E_n \subset \R$ be disjoint Borel sets (we assume $\sup E_j < \inf E_{j+1}$). Then,
\begin{equation}
	\frac{1}{\pZ{n}{}} \int_{E_1} \cdots \int_{E_n} \det[p_t(a_i, x_j)]_{i,j = 1}^{n} \dd x_1 \cdots \dd x_n
\end{equation}
is equal to the probability that each path $X_j$ belongs to the set $E_j$ at time $t$, without any intersection between paths in the time interval $[0, t]$, for some normalizing constant $Z_n$.
\end{thm}

However, this is not a determinantal process, since the correlation functions are not expressible in terms of the determinant of a kernel. If we additionally condition the paths to end at time $T > 0$ at some given points $b_1 < \cdots < b_n$, without any intersection between the paths along the whole time interval $[0, T ]$, then it can be shown (using an argument again based on the Markov property) that the random positions of the $n$ paths at a given time $t \in [0, T]$ have the joint probability density function
\begin{equation}
	\frac{1}{\pZ{n}{}} \det[p_t(a_i, x_j)]_{i,j = 1}^{n} \det[p_{T-t}(x_i, b_j)]_{i,j = 1}^{n} = \det[K_n(x_i, x_j)]_{i,j = 1}^{n}
	\label{Equation: biorthogonal pdf}
\end{equation} 
with a suitable normalizing constant $\pZ{n}{}$ and kernel
\begin{equation}
	K_n(x, y) \deff \sum_{j=1}^n \phi_j(x) \psi_j(y),
\end{equation}
\begin{equation}
	\phi_j \in \text{span}\{p_t(a_i, x)\}, \quad \psi_k \in \text{span}\{p_{T-t}(x, b_i)\},
\end{equation}
\begin{equation}
	\int_\R \phi_j(x) \psi_k(x) \dd x = \delta_{jk}.
\end{equation}


A interesting case is the confluent regime when two or more starting (or ending) points collapse together. For example, in the confluent limit as $a_j \rightarrow a$ and $b_j \rightarrow b$, for all $j$’s, applying L’Hopital rule to \eqref{Equation: biorthogonal pdf} gives
\begin{equation}
	\frac{1}{\tilde{Z}}  \det\left[\frac{\dd^{i-1}}{\dd a^{i-1}} p_t(a_i, x_j)\right]_{i,j = 1}^{n} \det\left[\frac{\dd^{i-1}}{\dd b^{i-1}} p_{T-t}(x_i, b_j)\right]_{i,j = 1}^{n} 
\end{equation}
which is still a determinantal point process with kernel derived along the same method as before.

\section{Limit and Transformations of Determinantal Point Processes}

Let $f \in C^1(\R)$ be a bijection , with continuously differentiable inverse $f^{-1} = g \in \C^1(\R)$. Then, $f$ maps configurations on $\R$ to configuration on $\R$. The push-forward of a point process $\mathbb{P}$ under such mapping is $f(\mathbb{P})$ and it is again a point process.

\begin{prop} \cite{ManuelaGirottiPhD}
	If K is the kernel of a DDP $\mathbb{P}$, then $f(\mathbb{P})$ is also a DDP with kernel
	\begin{equation}
		\hat{K}(x, y) = \sqrt{g'(x) g'(y)} K(g(x), g(y)).
	\end{equation}
	where $f^{-1} = g$.
\end{prop}

In particular, if we consider a linear function $f(x) = \alpha (x - x^*),$ for $\alpha > 0$ and $x^* \in \R$, then the transformed point process $f(\mathbb{P})$ is a scaled and translated version of the original point process $\mathbb{P}$. The new correlation kernel is
\begin{equation}
	\hat{K}(x, y) = \frac{1}{\alpha} K\left(x^* + \frac{x}{\alpha}, x^* + \frac{y}{\alpha}\right).
\end{equation}

Others transformations are also possible. Consider what happens when one fixes $x^* \in \R$ and, given that $x^* \in \chi$, removes $x^*$ from the configuration. The resulting process $\chi \setminus \{x^*\}$ is a new point process.
 
 \begin{prop} \cite{ManuelaGirottiPhD}
 		Let K be the correlation kernel of a DDP and let $x^*$ be a point such that $0 < K(x^*, x^*) < + \infty$. Then, the point process obtained by removing $x^*$ is determinantal with kernel
 		\begin{equation}
 			\hat{K}(x, y) = K(x, y) - \frac{K(x, x^*) K(x^*, x)}{K(x^*, x^*)}.
 		\end{equation}
 \end{prop}
 
 We now want to consider sequences of finite determinantal point processes $\mathbb{P}_n$ with correlation kernel $K_n$ and understand how the limit behaves. 
 
 \begin{prop} \cite{ManuelaGirottiPhD}
	Let $\mathbb{P}$ and $\mathbb{P}_n$ be determinantal point processes with kernels $K$ and $K_n$ respectively. Let $K_n$ converge pointwise to $K$
	\begin{equation}
		\lim_{n \rightarrow \infty} K_n(x, y) = K(x, y)
	\end{equation}
	uniformly in $x, y$ over compact subsets of $\R$. Then, the point processes $\mathbb{P}_n$ converge to $\mathbb{P}$ weakly. 
 \end{prop}

A very important case studied here is of the limit of an scaled and centered determinantal point process. Suppose a sequence of kernels $K_n$ and fixed point $x^*$ of interest. We first perform a centering and rescaling of the form
\begin{equation}
	x \mapsto Cn^\alpha (x - x^*)
\end{equation} 
for some suitable value of $C, \gamma > 0$. Then the scaled kernels have a limit
\begin{equation}
	\lim_{n \rightarrow \infty} \frac{1}{C n^\alpha} K_n \left( x^* + \frac{x}{C n^\alpha}, x^* + \frac{y}{C n^\alpha} \right) = K(x, y)
\end{equation}
Therefore, the scaling limit $K$ is a kernel that corresponds to a determinantal point process with
an infinite number of points.  Interestingly, in many situations, the same scaling limit $K$ may appear. The phenomenon in known as \textbf{universality} in RMT. Examples of such kernels are the sine kernel
\begin{equation}
	K_{\sin} = \frac{\sin \pi(x - y)}{x - y}
\end{equation}
and the Airy kernel
\begin{equation}
	K_{\Ai} = \frac{\Ai(x)\Ai'(y) - \Ai'(x) \Ai(y)}{x - y}
\end{equation}
where $\Ai$ is the Airy function.